
O rendimento do moitão é dado pela equação:

\begin{equation}
  \eta_{\text{moitão}} = \frac{1}{K}\frac{\left(1-{\eta_{\text{rolamento}}}^K\right)}{\left(1-\eta_{\text{rolamento}}\right)}
  \label{eq:rendimento_moitao}
\end{equation}

Onde:

\begin{itemize}[leftmargin=1.25cm]
  \item $\eta_{\text{moitão}}$: rendimento do moitão.
  \item $K$: meio-ramal.
  \item $\eta_{\text{rolamento}}$: rendimento do rolamento.
\end{itemize}

\vspace{1em}

O meio-ramal é dado pela equação:

\begin{equation}
  K = \frac{n_c}{2}
  \label{eq:meio_ramal}
\end{equation}

Onde:

\begin{itemize}[leftmargin=1.25cm]
  \item $K$: meio-ramal.
  \item $n_c$: número de cabos.
\end{itemize}

\vspace{1em}

Substituindo os valores na equação \eqref{eq:meio_ramal}:

\begin{align*}
  K & = \frac{4}{2} \\
  K & = \boxed{2}
\end{align*}

O rendimento do rolamento adotado é de $0{,}985$.

Substituindo os valores na equação \eqref{eq:rendimento_moitao}:

\begin{align*}
  \eta_{\text{moitão}} & = \frac{1}{2}\frac{\left(1-0{,}985^2\right)}{\left(1-0{,}985\right)} \\
  \eta_{\text{moitão}} & = \boxed{0{,}9925}
\end{align*}
